% =====================================================================
%  Lexikální analýza
% =====================================================================
\section{Lexikální analýza}\label{sec:lexer}

\subsubsection*{Popis automatu}
Lexer implementuje konečný automat viditelný na obrázcích~\ref{fig:ka-id}
až~\ref{fig:ka-retezce}, které mají společný počáteční stav \texttt{s}.
Stavy \texttt{q} nepřijímající, stavy \texttt{f} přijímající a stavy
\texttt{p} pro zpracování komentáře. Lexer čte nejdelší platný lexém. Když
skončí v~přijímajícím stavu, vrátí token a nepoužitý znak nechá pro další
volání. Neplatný přechod z~nepřijímajícího stavu je lexikální chyba
s~návratovým kódem 1.

\subsubsection*{Odchylka: Vnořené komentáře}
Blokové komentáře lze libovolně vnořovat, což konečný automat sám rozpoznat
nedokáže. Stav \texttt{p2} proto používá čítač \texttt{depth}: \texttt{/*} ho
zvýší, \texttt{*/} sníží a při nule komentář končí. Neuzavřený komentář je
lexikální chyba.

\subsubsection*{Odchylka: Dopředné čtení}
Lexer drží aktuální znak \texttt{current} a jeden znak dopředu,
\texttt{lookahead}. Díky němu rozliší desetinnou tečku od samostatného tokenu
(\texttt{1.5} oproti \texttt{1.foo}) a v~komentářích pozná dvojice
\texttt{/*} a \texttt{*/}.

\subsubsection*{Odchylka: Dodatečná validace}
Po přečtení identifikátoru rozhodne funkce \texttt{keyword()}, jestli jde
o~klíčové nebo rezervované slovo. U~číselných literálů lexer kontroluje
počáteční nuly a pomocí \texttt{strtod} rozsah typu \texttt{Double}. Hlídá
i~hranice tokenů, třeba když za číslem nebo \texttt{@main} pokračuje znak
identifikátoru. Původní zápis lexému se nemění.

\subsubsection*{Odchylka: Logické řádky a odsazení}
Bílé znaky, komentáře a odsazení zpracovává lexer mimo samotný automat.
\texttt{TOKEN\_NEWLINE} vrací jen na konci neprázdného logického řádku mimo
závorky. Kontroluje pokračování řádků, tabulátory v~odsazení a pravidla pro
víceřádkové komentáře. Tokeny \texttt{INDENT} a \texttt{DEDENT} nepoužívá,
parser dostane odsazení a příznak na začátku řádku přímo ve struktuře
\texttt{token\_t}.

% diagramy v TikZu: obrazky/ka-*.tex (všechny začínají ve stavu s)

\begin{figure}[htbp]
  \centering
  % KA lexeru -- identifikátory, klíčová slova, @main, samotné _
\begin{tikzpicture}[ka,px]
  \node[state,initial]   (s)   at (208,266)  {$s$};
  \node[state,accepting] (f31) at (458,147)  {$f_{31}$};
  \node[state,accepting] (f1)  at (818,266)  {$f_{1}$};
  \node[state]           (q2)  at (408,430)  {$q_{2}$};
  \node[state]           (q3)  at (628,430)  {$q_{3}$};
  \node[state]           (q4)  at (848,430)  {$q_{4}$};
  \node[state]           (q5)  at (1068,430) {$q_{5}$};
  \node[state,accepting] (f2)  at (1288,430) {$f_{2}$};

  \node[tok,below=0.5mm of f31]            {TOKEN\_UNDERSCORE};
  \node[tok,below=1mm of f1,align=center]  {TOKEN\_IDENTIFIER\\/ klíčové slovo};
  \node[tok,below=1mm of f2]               {TOKEN\_MAIN};

  \draw[hrana] (s) -- (320,266) |- node[lbl,above,pos=0.75] {\kaPod} (f31);
  \draw[hrana] (s) -- node[lbl,below,pos=0.5] {\kaPis} (f1);
  \draw[hrana] (f31) -| (712,266) -- (f1);
  \node[lbl,above right] at (500,147) {\kaPis\,+\,\kaCis\,+\,\kaPod};
  \path (f1) edge[loop right] node[right] {\kaPis\,+\,\kaCis\,+\,\kaPod} ();
  \draw[hrana] (s) |- node[lbl,above,pos=0.75] {\texttt{@}} (q2);
  \draw[hrana] (q2) -- node[lbl,above] {\texttt{m}} (q3);
  \draw[hrana] (q3) -- node[lbl,above] {\texttt{a}} (q4);
  \draw[hrana] (q4) -- node[lbl,above] {\texttt{i}} (q5);
  \draw[hrana] (q5) -- node[lbl,above] {\texttt{n}} (f2);
\end{tikzpicture}

  \caption{Konečný automat lexeru -- identifikátory, klíčová slova
           a \texttt{@main}}
  \label{fig:ka-id}
\end{figure}

\begin{figure}[htbp]
  \centering
  % KA lexeru -- komentáře (rozšíření o čítač depth)
\begin{tikzpicture}[ka,px]
  \node[state,initial]   (s)   at (88,300)  {$s$};
  \node[state,accepting] (f26) at (338,300) {$f_{26}$};
  \node[state]           (p1)  at (638,300) {$p_{1}$};
  \node[state]           (p2)  at (638,560) {$p_{2}$};

  \node[tok, below=1mm of f26] {TOKEN\_SLASH};
  \node[pozn,below=1mm of p1]  {řádkový komentář};
  \node[pozn,below=1mm of p2]  {blokový komentář};

  \draw[hrana] (s) -- node[lbl,above] {\texttt{/}} (f26);
  \draw[hrana] (f26.east) -- (462,300) |- node[lbl,above,pos=0.75] {\texttt{/}} (p1);
  \path (p1) edge[loop right] node[lbl,right=1pt] {\kaKrome{LF}} ();
  \draw[hrana] (p1) |- node[lbl,above,pos=0.75] {LF / zpracuj konec řádku} (88,155) -- (s);
  \draw[hrana] (f26.east) -- (462,300) |- (p2);
  \node[lbl,left] at (462,485) {\texttt{*} / $\mathit{depth}=1$};
  \path (p2) edge[loop above] node[lbl] {\texttt{/*} / $\mathit{depth}{+}{+}$} ();
  \path (p2) edge[loop right] node[lbl,right=2pt] {\texttt{*/} a $\mathit{depth}>1$ / $\mathit{depth}{-}{-}$} ();
  \draw[hrana] (p2.west) -- (492,560) |- (492,720) -| (s);
  \node[lbl,below] at (290,720) {\texttt{*/} a $\mathit{depth}=1$ / $\mathit{depth}{-}{-}$; dokonči komentář};
\end{tikzpicture}

  \caption{Konečný automat lexeru -- komentáře
           (rozšíření o~čítač $\mathit{depth}$)}
  \label{fig:ka-komentare}
\end{figure}

\begin{figure}[htbp]
  \centering
  % KA lexeru -- operátory a oddělovače
% oba uzly s = tentýž stav
\begin{tikzpicture}[ka,px,every state/.append style={minimum size=7.5mm}]
  % --- levá část: jednoznakové tokeny ---------------------------------
  \node[state,initial] (s) at (88,616) {$s$};
  \foreach \i/\id/\sym/\tk/\x in {
      0/3/(/LPAREN/180,  1/4/)/RPAREN/193,  2/5/:/COLON/207,
      3/6/{,}/COMMA/221, 4/7/;/SEMICOLON/235, 5/8/+/PLUS/249,
      6/9/-/MINUS/263,   7/10/*/STAR/277,   8/33/[/LBRACKET/291,
      9/34/]/RBRACKET/305, 10/35/\{/LBRACE/319, 11/36/\}/RBRACE/333} {
    \node[state,accepting] (f\id) at (448,236+110*\i) {$f_{\id}$};
    % popisky vpravo od stavů
    \node[tok,right=0.5mm of f\id] {TOKEN\_\tk};
    \draw[hrana] (s.east) -- (\x,616) |- (f\id);
    \node[lbl,above] at (362,236+110*\i) {\texttt{\sym}};
  }

  % --- pravá část: víceznakové operátory -----------------------------
  \node[state,initial] (t)   at (838,616)   {$s$};
  \node[state,accepting] (f11) at (1138,276)  {$f_{11}$};
  \node[state,accepting] (f14) at (1438,276)  {$f_{14}$};
  \node[state,accepting] (f12) at (1138,496)  {$f_{12}$};
  \node[state,accepting] (f15) at (1438,496)  {$f_{15}$};
  \node[state,accepting] (f30) at (1438,606)  {$f_{30}$};
  \node[state,accepting] (f13) at (1138,716)  {$f_{13}$};
  \node[state,accepting] (f16) at (1438,716)  {$f_{16}$};
  \node[state,accepting] (f27) at (1138,936)  {$f_{27}$};
  \node[state,accepting] (f17) at (1438,936)  {$f_{17}$};
  \node[state]           (q19) at (1138,1066) {$q_{19}$};
  \node[state,accepting] (f28) at (1438,1066) {$f_{28}$};
  \node[state,accepting] (f32) at (1138,1196) {$f_{32}$};
  \node[state,accepting] (f29) at (1438,1196) {$f_{29}$};

  \foreach \n/\tk in {f11/ASSIGN, f14/EQUAL, f12/LESS, f15/LESS\_EQUAL,
                      f30/LARROW, f13/GREATER, f16/GREATER\_EQUAL,
                      f27/NOT, f17/NOT\_EQUAL, f28/AND, f32/PIPE, f29/OR}
    \node[tok,below=0.2mm of \n] {TOKEN\_\tk};
  \node[pozn,below=0.2mm of q19,xshift=4mm] {samotné \& je chyba};

  \foreach \n/\sym/\x in {f11/=/930, f12/</944, f13/>/958, f27/!/972,
                          q19/\&/990, f32/|/1010}
    \draw[hrana] (t.east) -- (\x,616) |-
      node[lbl,above,pos=0.8] {\texttt{\sym}} (\n);
  \draw[hrana] (f11) -- node[lbl,above] {\texttt{=}} (f14);
  \draw[hrana] (f12) -- node[lbl,above] {\texttt{=}} (f15);
  \draw[hrana] (f12.east) -- (1252,496) |- node[lbl,above,pos=0.75] {\texttt{-}} (f30);
  \draw[hrana] (f13) -- node[lbl,above] {\texttt{=}} (f16);
  \draw[hrana] (f27) -- node[lbl,above] {\texttt{=}} (f17);
  \draw[hrana] (q19) -- node[lbl,above] {\texttt{\&}} (f28);
  \draw[hrana] (f32) -- node[lbl,above] {\texttt{|}} (f29);
\end{tikzpicture}

  \caption{Konečný automat lexeru -- operátory a oddělovače
           (oba uzly $s$ jsou tentýž počáteční stav)}
  \label{fig:ka-op}
\end{figure}

\begin{figure}[htbp]
  \centering
  % KA lexeru -- číselné literály a tečka
\begin{tikzpicture}[ka,px]
  \node[state,initial]   (s)   at (96,467)   {$s$};
  \node[state,accepting] (f22) at (96,277)   {$f_{22}$};
  \node[state,accepting] (f18) at (376,467)  {$f_{18}$};
  \node[state]           (q7)  at (376,706)  {$q_{7}$};
  \node[state]           (q8)  at (676,467)  {$q_{8}$};
  \node[state,accepting] (f19) at (976,467)  {$f_{19}$};
  \node[state]           (q9)  at (976,706)  {$q_{9}$};
  \node[state,accepting] (f21) at (1276,467) {$f_{21}$};
  \node[state]           (q12) at (1276,677) {$q_{12}$};
  \node[state]           (q10) at (676,906)  {$q_{10}$};
  \node[state]           (q11) at (976,906)  {$q_{11}$};
  \node[state,accepting] (f20) at (1276,906) {$f_{20}$};

  \node[tok,above=1mm of f22]  {TOKEN\_DOT};
  \node[tok,below=1mm of f18]  {TOKEN\_INTEGER};
  \node[tok,below=1mm of f19]  {TOKEN\_DOUBLE};
  \node[tok,below=1mm of f21]  {TOKEN\_DOUBLE};
  \node[tok,anchor=north west] at (f20.south east) {TOKEN\_DOUBLE};
  \node[pozn,below=1mm of q7,align=center]  {po \kaPod{} je\\nutná číslice};
  \node[pozn,below=1mm of q8,align=center]  {nutná desetinná\\číslice};
  \node[pozn,below=1mm of q9,align=center]  {po \kaPod{} je\\nutná číslice};
  \node[pozn,below=0.5mm of q12,align=center] {po \kaPod{} je\\nutná číslice};
  \node[pozn,anchor=north east,align=right] at (q10.south west) {za \texttt{e}/\texttt{E}\\musí být číslice};
  \node[pozn,below=1mm of q11,align=center] {po znaménku\\číslice};

  % s
  \draw[hrana] (s) -- node[lbl,above] {\kaCis} (f18);
  \draw[hrana] (s) -- node[lbl,left] {\large\texttt{.}} (f22);
  \draw[hrana] (f22) -| node[lbl,above,pos=0.25] {\kaCis} (f19);
  % celá část
  \path (f18) edge[loop above] node[lbl] {\kaCis} ();
  \draw[hrana] (f18.east) -- (492,467) |- (q7);
  \node[lbl,above] at (462,467) {\kaPod};
  \draw[hrana] (q7) -- (245,706) |- node[lbl,left,pos=0.25] {\kaCis} (f18);
  \path (q7) edge[loop above] node[lbl] {\kaPod} ();
  \draw[hrana] (f18) -- node[lbl,above,pos=0.8] {\large\texttt{.}} (q8);
  \draw[hrana] (f18.east) -- (530,467) |- (q10);
  \node[lbl,right] at (530,620) {\texttt{e}\,+\,\texttt{E}};
  % desetinná část
  \draw[hrana] (q8) -- node[lbl,above,pos=0.35] {\kaCis} (f19);
  \path (f19) edge[out=115,in=160,looseness=7] node[lbl,above] {\kaCis} (f19);
  \draw[hrana] (f19.east) -- (1092,467) |- (q9);
  \node[lbl,above] at (1050,467) {\kaPod};
  \draw[hrana] (q9) -- (845,706) |- node[lbl,left,pos=0.25] {\kaCis} (f19);
  \path (q9) edge[loop above] node[lbl] {\kaPod} ();
  \draw[hrana] (f19) -- node[lbl,above,pos=0.7] {\texttt{d}\,+\,\texttt{D}} (f21);
  \draw[hrana] (f19.east) -- (1110,467) |- (1110,828) -| (q10);
  \node[lbl,above] at (780,828) {\texttt{e}\,+\,\texttt{E}};
  % exponent
  \draw[hrana] (q10) -- node[lbl,above] {\texttt{+}\,+\,\texttt{-}} (q11);
  \draw[hrana] (q11) -- node[lbl,above] {\kaCis} (f20);
  \draw[hrana] (q10) |- (676,1080) -| node[lbl,above,pos=0.25] {\kaCis} (f20);
  \path (f20) edge[loop above] node[lbl] {\kaCis} ();
  \draw[hrana] (f20.west) -- (1160,906) |- (q12);
  \node[lbl,left] at (1160,800) {\kaPod};
  \draw[hrana] (q12) -- (1430,677) |- node[lbl,left,pos=0.25] {\kaCis} (f20);
  \path (q12) edge[loop above] node[lbl] {\kaPod} ();
  \draw[hrana] (f20.east) -- (1520,906) |- node[lbl,right,pos=0.25] {\texttt{d}\,+\,\texttt{D}} (f21);
\end{tikzpicture}

  \caption{Konečný automat lexeru -- číselné literály a tečka}
  \label{fig:ka-cisla}
\end{figure}

\begin{figure}[htbp]
  \centering
  % KA lexeru -- běžné a víceřádkové řetězce
\begin{tikzpicture}[ka,px]
  \node[state,initial]   (s)   at (96,377)   {$s$};
  \node[state]           (q13) at (346,377)  {$q_{13}$};
  \node[state,accepting] (f23) at (646,377)  {$f_{23}$};
  \node[state]           (q16) at (946,377)  {$q_{16}$};
  \node[state]           (q17) at (1196,377) {$q_{17}$};
  \node[state]           (q18) at (1446,377) {$q_{18}$};
  \node[state,accepting] (f25) at (1446,587) {$f_{25}$};
  \node[state]           (q14) at (346,687)  {$q_{14}$};
  \node[state]           (q15) at (646,687)  {$q_{15}$};
  \node[state,accepting] (f24) at (96,877)   {$f_{24}$};

  \node[pozn,below=1mm of q13]                 {první uvozovka};
  \node[tok, below=1mm of f23]                 {TOKEN\_STRING: \kaUv\kaUv};
  \node[pozn,anchor=north west,align=left] at (q16.south east) {raw: žádná\\závěrečná \kaUv};
  \node[pozn,anchor=north west,align=left] at (q17.south east) {raw: jedna\\závěrečná \kaUv};
  \node[pozn,anchor=south west,align=left] at (q18.north east) {raw: dvě\\závěrečné \kaUv};
  \node[tok, below=1mm of f25]                 {TOKEN\_STRING};
  \node[pozn,anchor=north west,align=left] at (q14.south east) {obsah běžného\\řetězce};
  \node[pozn,right=1mm of q15]                 {escape};
  \node[tok, below=1mm of f24]                 {TOKEN\_STRING};

  % běžný řetězec
  \draw[hrana] (s) -- node[lbl,above,pos=0.3] {\kaUv} (q13);
  \draw[hrana] (q13) -- node[lbl,above] {\kaUv} (f23);
  \draw[hrana] (q13.west) -- (228,377) |- (q14);
  \node[lbl,left,align=right] at (228,540) {\kaKrome{\kaUv\,+\,\kaBs\\+\,LF}};
  \path (q14) edge[loop above] ();
  \node[lbl,right] at (360,610) {\kaKrome{\kaUv\,+\,\kaBs\,+\,LF}};
  \draw[hrana] (q13.east) -- (470,377) |- (470,490) -| (q15);
  \node[lbl,below] at (560,490) {\kaBs};
  \draw[hrana] (q14) -- node[lbl,above] {\kaBs} (q15);
  \draw[hrana] (q15) |- (346,858) -- (q14);
  \node[lbl,above] at (496,858) {\kaBs\,+\,\kaUv\,+\,\texttt{n}\,+\,\texttt{r}\,+\,\texttt{t}};
  \draw[hrana] (q14.west) -- (96,687) -- node[lbl,left,pos=0.5] {\kaUv} (f24);
  % víceřádkový řetězec
  \draw[hrana] (f23) -- node[lbl,above] {\kaUv} (q16);
  \path (q16) edge[out=110,in=150,looseness=7] node[lbl,left] {\kaKrome{\kaUv}} (q16);
  \draw[hrana] (q16) -- node[lbl,above] {\kaUv} (q17);
  \draw[hrana] (q17) -- node[lbl,above] {\kaUv} (q18);
  \draw[hrana] (q17) |- (946,545) -- (q16);
  \node[lbl,below] at (1071,545) {\kaKrome{\kaUv}};
  \draw[hrana] (q18) |- (946,195) -- (q16);
  \node[lbl,above] at (1196,195) {\kaKrome{\kaUv}};
  \draw[hrana] (q18.east) -- (1560,377) |- node[lbl,right,pos=0.25] {\kaUv} (f25);
  \path (f25) edge[loop left] node[lbl] {\kaUv} ();
\end{tikzpicture}

  \caption{Konečný automat lexeru -- běžné a víceřádkové řetězce}
  \label{fig:ka-retezce}
\end{figure}
